\documentclass{lmcs}

% ── Personal macros ──
\theoremstyle{plain}
\newtheorem{satz}[thm]{Satz}

\theoremstyle{definition}
\newtheorem{notation}[thm]{Notation}

% ── Math ──
\newcommand{\NN}{\mathbb{N}}
\newcommand{\RR}{\mathbb{R}}
\newcommand{\id}{\mathsf{id}}
\newcommand{\Lam}[1]{\lambda #1.\,}
\newcommand{\app}[2]{#1\;#2}
\newcommand{\fix}{\mathsf{Fix}}
\newcommand{\quoteop}{\mathsf{Quote}}
\newcommand{\evalop}{\mathsf{Eval}}
\newcommand{\Tcal}{\mathcal{T}}
\newcommand{\Acal}{\mathcal{A}}
\newcommand{\Scal}{\mathcal{S}}
\newcommand{\etal}{\textit{et al.}}

\title[Fix as Primitive]{A Term Calculus with Fixed Points as Primitive:
  The T-Transform and the Unification of Operator Equations}

\author{Zixi Li}
\address{Sun Yat-sen University}
\email{lizx93@mail2.sysu.edu.cn}

\begin{document}

\begin{abstract}
We define a term calculus extending the simply-typed lambda-calculus with
three self-reference primitives---quotation, evaluation, and fixed
points---in which every operator equation on a vector space is solved by a
single construct: the fixed point of the T-transform, defined as
T(u) = u minus eta times the operator applied to u. We prove that among all
affine term-forming operations, the T-transform is the canonical form
satisfying the fixed-point-is-solution property. We give a decomposition into
differential and constraint operators that unifies scalar equations, initial
value problems, boundary value problems, and time-dependent PDEs as
instances of the same operator equation differing only in the choice of
constraint. We show that reverse-mode automatic differentiation is
expressible as a term-to-term transformation within the calculus, and that
for implicit-layer models, both the forward pass and the gradient are fixed
points of the T-transform. The entire system is implemented in approximately
1,500 lines of code. The solver kernel is under 30 lines.
\end{abstract}

\begin{keywords}
\mbox{Fixed-point} combinator, term calculus, operational semantics,
Kleene recursion theorem, operator equation, automatic differentiation,
deep equilibrium models, Green's function
\end{keywords}

\maketitle

% ═══════════════════════════════════════════════════════════════════════════
\section{Introduction}
% ═══════════════════════════════════════════════════════════════════════════

\subsection{Motivation}

Consider two programs:
\begin{itemize}
\item \texttt{f(x) = x + 1} --- a term that computes a value.
\item \texttt{solve(x**3 - 2*x - 5 = 0)} --- a call to a library that
  finds a value.
\end{itemize}
The first is a term in a programming language. The second is a black-box
procedure, invisible to the type system, opaque to composition. This
distinction is not a mathematical necessity; it is an engineering artifact.

We propose a term calculus in which solving an equation is not an operation
applied to a term, but \emph{a term itself}. The fixed-point combinator
$\fix$---not a numerical algorithm, but a constructor in the
syntax---constitutes the sole mechanism for solving equations. Given any
operator $A$ on a vector space and a scalar $\eta \neq 0$, the term
$\fix(\Lam{u} u - \eta \cdot A[u])$ \emph{is} the solution to $A[u] = 0$.
The proof that this term satisfies the equation is two lines and holds for
any $A$ on any space.

\subsection{The T-Transform}

The T-transform is a \emph{residual operator embedding}: rather than solving
$A[u] = 0$ directly, we embed the equation into a discrete dynamical system
whose stationary points are exactly the zeros of $A$. Define, for any operator
$A$ on a vector space $X$ and scalar $\eta \neq 0$:
\begin{equation}\label{eq:T}
  T_{A,\eta}[u] \equiv u - \eta \cdot A[u].
\end{equation}
This induces the operator flow $u_{k+1} = T_{A,\eta}(u_k)$, whose continuous
analog is $\partial_t u = -A[u]$ (explicit Euler discretization with time step
$\eta$). The central observation of this paper is that the fixed point of this
flow is exactly the solution:
\begin{equation}\label{eq:fix-solves}
  \fix(T) = T(\fix(T)) \;\Longrightarrow\; A[\fix(T)] = 0.
\end{equation}
The forward direction of~\eqref{eq:fix-solves} holds by the operational
semantics of $\fix$. The reverse direction---that $A[\fix(T)] = 0$
follows---is a two-line algebraic cancellation (Theorem~\ref{thm:fix-solves}).

We characterize the T-transform within the class of affine residual operators
(Theorem~\ref{thm:canonicity}): among all binary term-forming operations
$\Phi(u, v)$ affine in both arguments, the form
$\Phi(u, A[u]) = u - \eta \cdot A[u]$ is the canonical representative
satisfying the fixed-point-is-solution property. Any other linear combination
$\alpha u + \beta A[u]$ fails to entail $A[\fix(\Phi)] = 0$ from
$\fix(\Phi) = \Phi(\fix(\Phi))$ without additional hypotheses on $A$.

\subsection{Contributions}

\begin{enumerate}
\item \textbf{A term calculus with $\fix$ as primitive}
  (Section~\ref{sec:calculus}). We extend the $\lambda$-calculus with
  $\quoteop$, $\evalop$, and $\fix$ as syntactic constructors whose
  operational semantics directly encode self-reference. The evaluation rule
  $\fix(f) \to \app{f}{\fix(f)}$ implements Kleene's first recursion theorem
  as a reduction step.

\item \textbf{The T-transform and its canonicity}
  (Section~\ref{sec:ttransform}). We characterize the T-transform within the
  class of affine residual operators: among all linear combinations
  $\alpha u + \beta A[u]$, the fixed-point-is-solution property holds iff
  $\alpha = 1$ and $\beta \neq 0$. The T-transform is the canonical
  representative of this class; its form is forced by the requirement that
  $\fix(\Phi) = \Phi(\fix(\Phi))$ entail $A[\fix(\Phi)] = 0$ without
  additional hypotheses on $A$.

\item \textbf{Unification of equation types via $(D, C)$}
  (Section~\ref{sec:unified}). We show that scalar equations, IVPs, BVPs, and
  time-dependent PDEs differ only in the choice of the constraint
  operator~$C$ in the decomposition $(D, C)u = (f, g)$. The solver never
  branches on equation type.

\item \textbf{Automatic differentiation as a term operator}
  (Section~\ref{sec:ad}). We define a Wengert tape within the calculus and
  prove that reverse-mode AD is a term-to-term transformation. The gradient
  of any expression is computable as $\fix(T)$ applied to the adjoint
  equation.

\item \textbf{DEQ gradient as nested $\fix$}
  (Section~\ref{sec:deq}). We prove that for implicit-layer models, both the
  forward pass $z^* = \fix(\Lam{z} f(z, x, \theta))$ and the gradient
  $g^* = \fix(\Lam{g} \partial L/\partial z^* + g \cdot \partial f/\partial z)$
  are instances of $\fix(T)$.
\end{enumerate}

\subsection{Related Work}

Kleene~\cite{kleene1952} introduced the first recursion theorem as a
meta-theorem about general recursive functions; the second recursion theorem
appeared earlier in~\cite{kleene1938}.
Barendregt~\cite{barendregt1984} treats the Y combinator as a definable term
within the untyped $\lambda$-calculus. Our approach differs: $\fix$ is a
\emph{primitive} constructor, not definable from application and abstraction,
and its operational semantics \emph{is} Kleene's first recursion theorem.

Plotkin~\cite{plotkin1977} and Scott~\cite{scott1976} gave domain-theoretic
semantics for fixed-point operators in PCF and LCF. Our calculus is unityped:
all terms inhabit a single syntactic domain, without a separate type
hierarchy for functions and values.

The Wengert tape~\cite{wengert1964} is the standard data structure for
reverse-mode automatic differentiation. Its formulation as a term-to-term
transformation within a $\lambda$-calculus, rather than as an external
procedure, is new.

Bai \etal~\cite{bai2019} introduced Deep Equilibrium Models and the
implicit-function-theorem gradient. The use of fixed-point iteration for the
gradient computation was presented as a numerical technique. We recast both
forward and backward passes as instances of the same $\fix(T)$ construct, and
we prove that this recasting is not an analogy---it is a formal identity
within the calculus.

Chen \etal~\cite{chen2018} unified ODE solvers with residual networks. Their
adjoint method propagates gradients through a continuous-time dynamical
system. Our $\fix(T)$ framework subsumes their forward pass but replaces the
adjoint ODE with another fixed point, eliminating the need to integrate
backward in time.

% ═══════════════════════════════════════════════════════════════════════════
\section{The Term Calculus}\label{sec:calculus}
% ═══════════════════════════════════════════════════════════════════════════

\subsection{Syntax}

Let $\mathcal{V}$ be a countable set of variable names. The set $\Tcal$ of
terms is defined inductively:

\begin{defi}[Terms]
\begin{align*}
t \;::=&\;
  \mathsf{Var}(x)                     &&\text{variable, } x \in \mathcal{V}  \\
\mid&\; \mathsf{Lam}(x, t)            &&\text{abstraction, } \lambda x.\, t  \\
\mid&\; \mathsf{App}(t_1, t_2)        &&\text{application}                    \\
\mid&\; \mathsf{Quote}(t)             &&\text{reification}                     \\
\mid&\; \mathsf{Eval}(t_1, t_2)       &&\text{self-application}                \\
\mid&\; \fix(t)                       &&\text{fixed point}                     \\
\mid&\; \mathsf{Nat}(n)               &&\text{natural number, } n \in \NN      \\
\mid&\; \mathsf{Cons}(t_1, t_2)       &&\text{pair}                            \\
\mid&\; \mathsf{Fun}(f, S)            &&\text{quantity in space } S            \\
\mid&\; \mathsf{Space}(\mathit{name}, D, \delta) &&\text{vector space}         \\
\mid&\; \mathsf{Integ}(t, a)          &&\text{integral operator (Volterra)}    \\
\mid&\; \mathsf{Integ}(t, a, b)       &&\text{integral operator (Fredholm)}    \\
\mid&\; \mathsf{Diff}(t)              &&\text{derivative operator}
\end{align*}
\end{defi}

The constructors fall into four groups:
\begin{itemize}
\item \textbf{Computation}: $\mathsf{Var}$, $\mathsf{Lam}$, $\mathsf{App}$
  --- the pure $\lambda$-calculus.
\item \textbf{Self-reference}: $\mathsf{Quote}$, $\mathsf{Eval}$, $\fix$
  --- reification, self-application, and fixed points.
\item \textbf{Data and quantities}: $\mathsf{Nat}$, $\mathsf{Cons}$,
  $\mathsf{Fun}$, $\mathsf{Space}$ --- natural numbers, pairs, vector-space
  points, and the spaces themselves.
\item \textbf{Analysis}: $\mathsf{Integ}$, $\mathsf{Diff}$ --- integral and
  derivative as first-class terms.
\end{itemize}

There is no primitive type distinction between ``functions'' and ``values''.
$\mathsf{Fun}(f, S)$ is a term like any other. A scalar $3.14$ is
$\mathsf{Fun}(\Lam{x} 3.14, \mathsf{SPACE\_R})$---a constant function on the
real line. This unification is essential: it ensures that the T-transform
$T_{A,\eta}[u] = u - \eta \cdot A[u]$ is well-typed for any operator $A$,
regardless of whether $u$ is a scalar (optimization), a function (PDE
solving), or a parameter vector (training).

\subsection{Operational Semantics}

The reduction relation $\to$ is defined by the following rules. Evaluation
proceeds to weak head normal form.

\begin{defi}[Reduction rules]
\begin{description}
\item[$\beta$-reduction:]
  $\mathsf{App}(\mathsf{Lam}(x, t), s) \to t[x := s]$

\item[Self-application:]
  $\evalop(\quoteop(t), s) \to \app{t}{s}$

\item[Fixed point:]
  $\fix(f) \to \app{f}{\fix(f)}$

\item[Integral application (Volterra):]
  $\mathsf{App}(\mathsf{Integ}(\mathsf{Lam}(\tau, b), a), x)
   \to F(x)$ \\
  where $F(x) = \int_a^x b[\tau := \xi]\,\mathrm{d}\xi$

\item[Integral application (Fredholm):]
  $\mathsf{App}(\mathsf{Integ}(\mathsf{Lam}(\tau, b), a_0, a_1), x)
   \to G(x)$ \\
  where $G(x) = \int_{a_0}^{a_1} K(x, \xi) \cdot b[\tau := \xi]\,
   \mathrm{d}\xi$

\item[Derivative application:]
  $\mathsf{App}(\mathsf{Diff}(\mathsf{Lam}(x, b)), p)
   \to D(b)[x := p]$ \\
  where $D(b) = \lim_{h \to 0} \frac{b[x := p+h] - b[x := p-h]}{2h}$

\item[Fun evaluation:]
  $\mathsf{App}(\mathsf{Fun}(f, S), x) \to f(x)$
\end{description}
\end{defi}

The recursion depth of $\fix$ expansion is bounded by a global parameter
$\mathit{MAX\_FIX}$ to ensure termination of evaluation. The formal
correctness of fixed-point constructions does not depend on this bound;
it is a practical measure for the interpreter.

\begin{rem}[Semantic interpretation of $\fix$]\label{rem:fix-semantics}
The reduction rule $\fix(f) \to \app{f}{\fix(f)}$ defines the
\emph{operational} semantics of $\fix$. For the theorems of
Sections~\ref{sec:ttransform}--\ref{sec:unified}, the \emph{denotational}
interpretation of $\fix$ is the Banach fixed point in a complete metric space:
when $X$ is a complete metric space and $T : X \to X$ is a contraction, $\fix(T)$
denotes the unique element $u^* \in X$ satisfying $u^* = T(u^*)$
(Theorem~\ref{thm:contractivity}). The operational rule and the denotational
interpretation are compatible: the reduction sequence $u_{k+1} = T(u_k)$
converges to the denotational fixed point when $T$ is a contraction. The
syntactic theorem $\fix(T) = T(\fix(T))$ (Theorem~\ref{thm:fix-solves}) holds
in both interpretations---as a reduction step in the operational semantics, and
as the defining equation of the Banach fixed point in the denotational
semantics. We do not mix domain-theoretic, rewrite, and metric interpretations:
throughout this paper, the denotational semantics of $\fix$ is the
Banach fixed point over complete metric spaces.
\end{rem}

\subsection{The Environment and Numerical Backend}

Evaluation occurs in an environment $\mathit{Env}$ that maps variable names
to values. The core evaluator $\mathsf{eval\_term}(t, \mathit{env}, fd)$
evaluates $t$ to weak head normal form via repeated reduction. The
parameter $fd$ tracks the current $\fix$ depth to enforce the
$\mathit{MAX\_FIX}$ bound.

\begin{defi}[Environment]
$\mathit{Env}$ is a chained lookup structure: each frame maps a finite set of
variable names to terms, with a parent pointer to the enclosing frame. The
lookup function $\mathsf{lookup}(\mathit{env}, x)$ searches frames from
innermost to outermost.
\end{defi}

The evaluation loop handles $\fix$ specially. When the evaluator encounters
$\fix(f)$, it applies one step of expansion:
\[
\mathsf{eval\_term}(\fix(f)) \to \mathsf{eval\_term}(\mathsf{App}(f, \fix(f)), fd+1).
\]
If $fd$ exceeds $\mathit{MAX\_FIX}$, evaluation terminates with the current
approximation. This is a pragmatic bound---the formal correctness of
$\fix(T)$ as a solution (Theorem~\ref{thm:fix-solves}) holds regardless of
how many expansion steps are taken, because the identity
$\fix(T) = T(\fix(T))$ is an equation in the syntax, not an asymptotic
limit. The $\mathit{MAX\_FIX}$ parameter is set to 5000 in the implementation,
which is far larger than any fixed-point iteration that converges in
practice.

\subsection{Arithmetic and the Critical Compilation Rule}

The calculus lifts arithmetic pointwise over function spaces. Binary
operations ($+$, $-$, $\times$) are defined via a single dispatcher
$\mathsf{\_binop}$ that handles all combinations:

\begin{defi}[Pointwise arithmetic]
For terms $f, g$ in the same space, $\mathsf{\_binop}(f, g, op)$
computes $h(x) = op(f(x), g(x))$ pointwise:
\begin{enumerate}
\item $(\mathsf{Fun}, \mathsf{Fun}) \to \mathsf{Fun}$: compose the callables
  pointwise.
\item $(\mathsf{Fun}, \mathsf{Lam}) \to \mathsf{Fun}$: \textbf{compile
  $\mathsf{Lam}$ to $\mathsf{Fun}$ first}, then combine as
  $(\mathsf{Fun}, \mathsf{Fun})$.
\item $(\mathsf{Lam}, \mathsf{Lam})$ with the same parameter: push the
  operation symbolically into the body (for construction, not iteration).
\end{enumerate}
\end{defi}

The second rule is critical. A na\"ive implementation would push the
$\mathsf{Fun}$ into the $\mathsf{Lam}$ body:
$\mathsf{Lam}(p, \mathsf{App}(\dots, p))$. After $n$ iterations, the
$\mathsf{Lam}$ nesting depth is $O(n)$, and evaluation traverses $n$ nested
closures. The compile rule avoids this entirely: $\mathsf{\_compile\_lam}$
samples the $\mathsf{Lam}$ on a grid spanning the function's domain in one
shot, returning a flat $\mathsf{Fun}$ backed by piecewise linear
interpolation, which never triggers nested $\mathsf{eval\_term}$ calls. This
is the single design rule that makes the fixed-point iteration practical.

After each Euler step, $\mathsf{flatten\_term}$ samples the resulting
$\mathsf{Fun}$ on a fresh grid and wraps it in a new interpolation, breaking
the closure chain. Without flattening, the $n$-th iterate chains through $n$
closures; with flattening, each iterate is a flat lookup table whose
evaluation cost is $O(1)$ regardless of iteration count.

\subsection{Free Variables and Substitution}

\begin{defi}[Free variables]
The set of free variables $\mathsf{FV}(t)$ is defined by structural
induction:
\begin{align*}
\mathsf{FV}(\mathsf{Var}(x)) &= \{x\} \\
\mathsf{FV}(\mathsf{Lam}(x, t)) &= \mathsf{FV}(t) \setminus \{x\} \\
\mathsf{FV}(\mathsf{App}(t_1, t_2)) &= \mathsf{FV}(t_1) \cup
  \mathsf{FV}(t_2) \\
\mathsf{FV}(\fix(t)) &= \mathsf{FV}(t) \\
\mathsf{FV}(\mathsf{Quote}(t)) &= \mathsf{FV}(t) \\
\mathsf{FV}(\mathsf{Eval}(t_1, t_2)) &= \mathsf{FV}(t_1) \cup
  \mathsf{FV}(t_2) \\
\mathsf{FV}(\mathsf{Integ}(t, a)) &= \mathsf{FV}(t) \cup \mathsf{FV}(a) \\
\mathsf{FV}(\mathsf{Integ}(t, a, b)) &= \mathsf{FV}(t) \cup \mathsf{FV}(a)
  \cup \mathsf{FV}(b) \\
\mathsf{FV}(\mathsf{Diff}(t)) &= \mathsf{FV}(t)
\end{align*}
\end{defi}

Substitution $t[x := s]$ is capture-avoiding and defined for all term types.
The only binding construct is $\mathsf{Lam}$;
$\mathsf{Integ}(\mathsf{Lam}(\tau, b), \dots)$ binds $\tau$ in $b$ only
during integration---the binder is structural, not a separate binding form.

\subsection{Fix Implements Kleene}

\begin{thm}[Fix implements Kleene]\label{thm:kleene}
The reduction rule $\fix(f) \to \app{f}{\fix(f)}$ directly implements
Kleene's first recursion theorem: for any computable function $f$, there
exists a term $\fix(f)$ such that $\fix(f)$ reduces to $f(\fix(f))$.
\end{thm}

\begin{proof}
The reduction rule IS the statement. The existence of the fixed point is
not a meta-theorem about the calculus; it is a reduction step within the
calculus. $\fix$ is the witness whose self-application generates the
Kleene fixed point by construction.
\end{proof}

\begin{cor}[Self-reference completeness]\label{cor:selfref}
The triple $(\quoteop, \evalop, \fix)$ is complete for self-referential
constructions: any computable function of its own description can be
expressed.
\end{cor}

\begin{proof}
$\quoteop$ reifies a term as data accessible to other terms. $\evalop$
applies a reified term to an argument, implementing the universal
function. $\fix$ constructs the self-application point. Together these
three operations suffice to implement any interpreter, self-modifying
program, or meta-circular evaluator~\cite{barendregt1984}.
\end{proof}

\subsection{Space Polymorphism}

A $\mathsf{Space}$ term carries its own algebraic structure:

\begin{defi}[Space]
\[
\mathsf{Space}(\mathit{name}, D, \delta)
\]
where $D \subseteq \RR^n$ is the domain and $\delta : \Tcal \times \Tcal \to
\RR^+$ is the distance function.
\end{defi}

The distance function $\delta$ is polymorphic: for finite-dimensional spaces
$\delta(a, b) = |a - b|$ (Euclidean), for function spaces $\delta(f, g) =
(\int |f - g|^2)^{1/2}$ ($L^2$).

\begin{lem}[Distance polymorphism]\label{lem:space-poly}
The solver evaluates $\mathsf{sp}.\mathsf{distance}(a, b)$ without
branching on the space type. The metric is a property of the space, not of
the solver.
\end{lem}

\begin{proof}
Each $\mathsf{Space}$ instance carries a closure implementing $\delta$. The
solver calls this closure uniformly. The dispatch to Euclidean or $L^2$
metric occurs inside the closure, invisible to the solver's control flow.
\end{proof}

The importance of space polymorphism for the T-transform framework cannot be
overstated. The T-transform $T_{A,\eta}[u] = u - \eta \cdot A[u]$ involves
subtraction---the operation $u - \eta \cdot A[u]$ must be meaningful for
whatever space $u$ inhabits. When solving a scalar equation, $u \in \RR$ and
subtraction is the usual one. When solving a BVP, $u$ is a function on $[0,1]$
and subtraction must be pointwise. When training a neural network, $u$ is a
parameter vector and $\eta \cdot A[u]$ is a gradient step in $\RR^p$. The
solver---the single function $\mathsf{solve}(A, \eta)$ that returns
$\fix(T(A, \eta))$---calls the same arithmetic operations regardless. Space
polymorphism makes this possible: the space object dispatches the arithmetic,
not the solver.

% ═══════════════════════════════════════════════════════════════════════════
\section{The T-Transform}\label{sec:ttransform}
% ═══════════════════════════════════════════════════════════════════════════

\subsection{Definition and Core Theorem}

The T-transform is a \emph{residual operator embedding}: it reinterprets the
static operator equation $A[u] = 0$ as the fixed-point condition of a discrete
dynamical system on the function space $X$. The central observation is that
\emph{solving an equation} and \emph{finding a stationary point of an
operator flow} are the same mathematical act.

Given an operator $A : X \to X$, the T-transform induces the one-parameter
family of discrete flows:
\[
u_{k+1} = T_{A,\eta}(u_k), \qquad
T_{A,\eta}[u] \equiv u - \eta \cdot A[u].
\]
The continuous-time analog is the operator differential equation
$\partial_t u = -A[u]$, of which $T_{A,\eta}$ is the explicit Euler
discretization with time step $\eta$. Every fixed point of the flow is a zero
of $A$, and every zero of $A$ is a fixed point of the flow---the correspondence
is exact (Theorem~\ref{thm:fix-solves}). What is unified across equation types
is not the structure of the equation but the \emph{evolution in function space}
that solves it.

\begin{defi}[T-transform]
For any operator $A$ on a vector space $X$ and any scalar $\eta \neq 0$:
\[
T_{A,\eta}[u] \equiv u - \eta \cdot A[u].
\]
In the term calculus, $T_{A,\eta}$ is the abstraction
$\mathsf{Lam}(u,\; \mathsf{sub}(u,\; \mathsf{mul}(\mathsf{R}(\eta),\;
\mathsf{App}(A, u))))$.
\end{defi}

\begin{thm}[Fixed-point is solution]\label{thm:fix-solves}
For any operator $A$ and any $\eta \neq 0$:
\[
\fix(T_{A,\eta}) = T_{A,\eta}(\fix(T_{A,\eta}))
\;\Longrightarrow\;
A[\fix(T_{A,\eta})] = 0.
\]
\end{thm}

\begin{proof}
\begin{align*}
\fix(T) &= T(\fix(T))
  &&\text{(definition of $\fix$)} \\
        &= \fix(T) - \eta \cdot A[\fix(T)]
  &&\text{(definition of $T$)} \\
\Rightarrow\; 0 &= -\eta \cdot A[\fix(T)]
  &&\text{(subtract $\fix(T)$)} \\
\Rightarrow\; A[\fix(T)] &= 0
  &&\text{($\eta \neq 0$)}
\end{align*}
\end{proof}

The proof is two algebraic steps. It applies to \emph{any} operator~$A$ on
\emph{any} vector space, without assumptions about linearity, continuity, or
differentiability. The term $\fix(T_{A,\eta})$ is a \emph{syntactic} theorem
that $A[u] = 0$ has a solution---the existence of the solution is guaranteed
by the reduction rule $\fix(f) \to \app{f}{\fix(f)}$.

\begin{rem}[Strength of the embedding]\label{rem:embedding-strength}
Theorem~\ref{thm:fix-solves} establishes a \emph{dynamical-system embedding}:
the equation $A[u] = 0$ is equivalent to the fixed-point condition
$u = T_{A,\eta}(u)$. This is not a claim of structural isomorphism across
equation classes (which would require categorical or operator-theoretic
equivalence proofs), but a claim of \emph{uniform representability}: every
operator equation admits the same syntactic form $\fix(T_{A,\eta})$, and the
semantic content resides entirely in the choice of $A$. The unification is at
the level of the solving process (evolution in function space), not at the
level of the equation structure.
\end{rem}

\subsection{Canonicity}

\begin{thm}[Canonicity of the T-transform]\label{thm:canonicity}
Among all term-forming operations $\Phi(u, v)$ affine in both arguments, the
form $\Phi(u, A[u]) = u - \eta \cdot A[u]$ is the canonical representative
satisfying the fixed-point-is-solution property
$\fix(\Phi) = \Phi(\fix(\Phi)) \Rightarrow A[\fix(\Phi)] = 0$.
\end{thm}

\begin{proof}
Let $\Phi(u, v) = \alpha u + \beta v$ with $\alpha, \beta$ scalars,
representing the most general linear term-forming operation on the pair
$(u, A[u])$. Then:
\begin{align*}
\fix(\Phi) &= \Phi(\fix(\Phi)) \\
           &= \alpha \cdot \fix(\Phi) + \beta \cdot A[\fix(\Phi)] \\[4pt]
\Rightarrow\; (1 - \alpha) \cdot \fix(\Phi) &= \beta \cdot A[\fix(\Phi)]
\end{align*}

For the solution condition $A[\fix(\Phi)] = 0$ to follow from
$\fix(\Phi) = \Phi(\fix(\Phi))$ without additional assumptions on $A$ or
$\fix(\Phi)$, we require:
\begin{enumerate}
\item $\beta \neq 0$: otherwise the equation gives no information about $A$.
\item $\alpha = 1$: otherwise $\fix(\Phi) = \frac{\beta}{1-\alpha} \cdot
  A[\fix(\Phi)]$ relates $\fix(\Phi)$ to $A[\fix(\Phi)]$ rather than forcing
  $A[\fix(\Phi)] = 0$.
\end{enumerate}
With $\alpha = 1$ and $\beta \neq 0$:
\[
0 \cdot \fix(\Phi) = \beta \cdot A[\fix(\Phi)]
\;\Longrightarrow\; A[\fix(\Phi)] = 0.
\]
Setting $\beta = -\eta$ gives the T-transform. \hfill $\square$
\end{proof}

\begin{cor}
The T-transform is canonical up to the choice of $\eta$. Different choices
of $\eta$ correspond to different preconditioners of the same dynamical
system $u_{t+1} = u_t - \eta \cdot A[u_t]$.
\end{cor}

\begin{rem}[Scope of the characterization]\label{rem:canonicity-scope}
Theorem~\ref{thm:canonicity} characterizes the T-transform within the
restricted hypothesis class of \emph{affine} term-forming operations
$\Phi(u, A[u]) = \alpha u + \beta A[u]$. It does not preclude the existence of
nonlinear preconditioners $\Phi(u, A[u])$ (involving, e.g., coordinate
transforms $\varphi(u)$ or state augmentation with $A^2[u], A^3[u], \dots$)
that also satisfy the fixed-point-is-solution property. The theorem establishes
that \emph{within the affine class}, the T-transform is the canonical and
essentially unique representative; extending the characterization to broader
operator classes is a direction for future work.
\end{rem}

The significance of Theorem~\ref{thm:canonicity} is that the T-transform is
not an arbitrary design choice. It is the \emph{canonical} linear combination
$u - \eta \cdot A[u]$ that makes the fixed-point-is-solution property hold
without additional hypotheses on $A$.

\subsection{Convergence}

\begin{thm}[Contractivity]\label{thm:contractivity}
Let $X$ be a complete metric space with distance $\delta$, and let
$A : X \to X$ be Lipschitz with constant $L$. If $|\eta| < 2/L$, then
$T_{A,\eta}$ is a contraction and $\fix(T_{A,\eta})$ is the unique fixed
point.
\end{thm}

\begin{proof}
For any $u, v \in X$:
\begin{align*}
\delta(T(u), T(v))
  &= \delta(u - \eta A[u],\; v - \eta A[v]) \\
  &= \| (I - \eta A)(u - v) \| \\
  &\leq \|I - \eta A\| \cdot \|u - v\|.
\end{align*}
For $\|I - \eta A\| < 1$, we require $|\eta| < 2/L$ when $A$ satisfies the
Baillon--Haddad condition $\langle Au - Av, u - v \rangle \geq
\frac{1}{L} \|Au - Av\|^2$. Under this condition $T$ is a strict contraction.
By the Banach fixed-point theorem~\cite{banach1922}, $\fix(T)$ exists and is
unique.
\end{proof}

\begin{rem}
The formal correctness of $\fix(T)$ (Theorem~\ref{thm:fix-solves}) does
\emph{not} depend on convergence. The fixed point \emph{is} the solution as
a syntactic theorem. Whether a particular numerical evaluation
converges---that is, whether the sequence $u_{k+1} = T(u_k)$ reaches a
neighborhood of $\fix(T)$ within finitely many iterations---is a separate
property of $A$ and $\eta$, governed by Theorem~\ref{thm:contractivity}.
This separation of syntactic correctness from numerical convergence is a key
architectural feature of the calculus.
\end{rem}

\subsection{The \texorpdfstring{$\eta$}{eta} Parameter}

\begin{thm}[$\eta$ unification]\label{thm:eta}
The scalar $\eta$ in $T[u] = u - \eta \cdot A[u]$ simultaneously plays three
roles:
\begin{enumerate}
\item \textbf{Time step} in differential equation solvers ($\eta = \Delta t$);
\item \textbf{Learning rate} in gradient-based optimization ($\eta =
  \mathrm{lr}$);
\item \textbf{Damping factor} in fixed-point iteration.
\end{enumerate}
These are not analogies. They are the same parameter in the same equation.
\end{thm}

\begin{proof}
In each case, $\eta$ controls the step size in the dynamics
$u_{t+1} = u_t - \eta \cdot A[u_t]$.
\begin{itemize}
\item For $A = d/dt$ (time derivative), $\eta = \Delta t$ recovers Euler's
  method.
\item For $A = \nabla_\theta \mathcal{L}$ (gradient), $\eta$ is the
  learning rate in gradient descent.
\item For $A = I - f$ (fixed-point residual), $\eta$ damps the iteration
  $u_{t+1} = (1-\eta)u_t + \eta f(u_t)$.
\end{itemize}
The optimal $\eta$ satisfies the same Armijo condition
$\mathcal{L}(u_t - \eta A[u_t]) \leq \mathcal{L}(u_t) - c \cdot \eta \cdot
\|A[u_t]\|^2$ regardless of whether $\mathcal{L}$ is a physical energy, a
loss function, or a residual norm. A schedule $\eta_t$ (line search, cosine
annealing) improves convergence identically across all three contexts.
\end{proof}

The unification of $\eta$ is visible directly in the implementation. The
training algorithm updates parameters via two lines:
\begin{center}
$\delta \leftarrow \eta \cdot \nabla_\theta \mathcal{L}$ \\
$\theta \leftarrow \theta - \delta$
\end{center}
This is exactly the T-transform $\theta \leftarrow \theta - \eta \cdot
\nabla_\theta \mathcal{L}$, with $\eta$ playing the role of learning rate.
The same $\eta$ appears in the solver construction:
\[
T_{A,\eta} \;=\; \Lam{u} u - \eta \cdot A[u],
\qquad
\mathsf{solve}(A, \eta) \;=\; \fix(T_{A,\eta}).
\]
Both the PDE solver and the neural network trainer invoke the same
$\fix(T)$ construct. The trainer is not analogous to the solver; the trainer
\emph{is} the solver, with $A = \nabla_\theta \mathcal{L}$.

% ═══════════════════════════════════════════════════════════════════════════
\section{The Unified Operator Equation}\label{sec:unified}
% ═══════════════════════════════════════════════════════════════════════════

\subsection{The \texorpdfstring{$(D, C)$}{(D, C)} Decomposition}

Every classical equation type decomposes into an interior law $D$ and an
exterior constraint $C$:

\begin{defi}[Unified operator equation]
\[
(D, C) u = (f, g)
\]
where
\begin{itemize}
\item $D : X \to Y$ is the differential operator (interior law);
\item $C : X \to Z$ is the constraint operator (boundary / initial slice);
\item $f \in Y$, $g \in Z$ are the data.
\end{itemize}
\end{defi}

The constraint operator $C$ selects a slice of the unknown function. The
equation type is determined entirely by the signature of $C$:

\begin{center}
\begin{tabular}{l l l}
\toprule
$C$ takes & Equation type & Canonical example \\
\midrule
$\emptyset$ & Scalar equation & $x^3 - 2x - 5 = 0$ \\
$u(a)$ & Initial value problem & $y' = y,\; y(0) = 1$ \\
$(u(a), u(b))$ & Boundary value problem & $-u'' = 1,\; u(0) = u(1) = 0$ \\
$(u(a), u(b)) \times \NN_t$ & Time-dependent PDE & $u_t = u_{xx},\;
  u(x,0) = \sin(\pi x)$ \\
\bottomrule
\end{tabular}
\end{center}

The operator $A$ in $A[u] = 0$ is then:
\[
A[u] = \begin{cases}
f(u) & C = \emptyset \quad\text{(scalar)} \\
(Du - f,\; Cu - g) & \text{otherwise}
\end{cases}
\]

\begin{lem}[Operator factorization]\label{lem:op-factor}
Let $X$ be a normed linear space. Every well-posed linear boundary value
problem induces a bounded linear operator factorization
$L = D_0 + BC : X \to Y \times Z$, where:
\begin{itemize}
\item $D_0 : X \to Y$ is the differential operator restricted to the kernel of
  $C$ (homogeneous constraints);
\item $C : X \to Z$ is the constraint (trace) operator selecting boundary or
  initial data;
\item $B : Z \to Y$ is the boundary lifting operator satisfying $CB = I_Z$ and
  $D_0B = 0$.
\end{itemize}
The operator $L$ is invertible iff the homogeneous problem $D_0 u = 0, Cu = 0$
has only the trivial solution $u = 0$.
\end{lem}

\begin{proof}
For any $u \in X$, write $u = u_0 + Bg$ where $Cu_0 = 0$ and $g = Cu$. Then
$Lu = (Du, Cu) = (D_0 u_0 + D B g, g)$. Since $D B$ maps $Z \to Y$, define the
Schur complement $S = D_0|_{\ker C}$ (invertible by well-posedness). The inverse
is $L^{-1}(f, g) = S^{-1}(f - D B g) + B g$, which is the Green's
representation of Theorem~\ref{thm:green}.
\end{proof}

This factorization makes precise the intuition that the constraint operator $C$
and the differential operator $D$ are not independent components but factors
of a single bounded linear operator on $X$. The $(D, C)$ notation is not merely
a pair of operators; it denotes the product structure $D_0 + BC$ in the
category of normed linear spaces.

\subsection{Green's Representation}

\begin{thm}[Green's representation]\label{thm:green}
For a linear operator $L = (D, C)$, the inverse $L^{-1}$ decomposes as:
\[
L^{-1}[(f, g)](x) = H[g](x) + \int_\Omega K(x, \xi) \cdot f(\xi)\,
\mathrm{d}\xi
\]
where $H$ propagates constraints ($Du = 0, Cu = g$) and $K$ is the Green's
kernel ($Du = \delta, Cu = 0$).
\end{thm}

\begin{proof}
By linearity of $L$. The homogeneous solution $H[g]$ satisfies
$L[H[g]] = (0, g)$. The particular solution with Green's kernel satisfies
$L[\int K f] = (f, 0)$. By superposition,
$L[H[g] + \int K f] = (f, g)$.
\end{proof}

The kernel type is determined by the order and domain of $D$:
\begin{itemize}
\item $D = d/dx$ on $[a, b]$: Volterra kernel (step function, causal). The
  integral equation $u(x) = u(a) + \int_a^x f(t)\,\mathrm{d}t$ propagates
  information forward---the value at $x$ depends only on $[a, x]$.
\item $D = -d^2/dx^2$ on $[a, b]$: Fredholm kernel (piecewise linear,
  symmetric). The Green's function $G(x, \xi) = \min(x, \xi) \cdot (1 -
  \max(x, \xi))$ couples every point to every other point---the value at $x$
  depends on the entire domain.
\end{itemize}
This distinction is not encoded in the solver. It is encoded in the arity of
$\mathsf{Integ}$: two arguments dispatch to Volterra, three to Fredholm. The
same term constructor, $\mathsf{Integ}$, handles both. Adding a fourth
argument would dispatch to a boundary-element kernel; adding a noise term
would dispatch to a stochastic integral. The constructor generalizes without
modifying the solver.

\subsection{Integ and Diff as First-Class Terms}

A single $\mathsf{Integ}$ constructor dispatches on the presence of the
third argument:

\begin{defi}[Integ dispatch]
\[
\mathsf{Integ}(f, a) \to \text{Volterra: } F(x) = \int_a^x f(t)\,\mathrm{d}t
\]
\[
\mathsf{Integ}(f, a, b) \to \text{Fredholm: }
F(x) = \int_a^b G(x, t) \cdot f(t)\,\mathrm{d}t
\]
\end{defi}

The evaluator dispatches on the arity of $\mathsf{Integ}$. No branching on
``what kind of equation is this'' occurs: the structure of $A$ determines the
evaluation strategy.

Similarly, $\mathsf{Diff}(\mathsf{Lam}(x, b))$ is a term of the calculus,
not a primitive escape to a numerical procedure. When applied to a point
$p$, it evaluates the symmetric difference quotient.

\subsection{Heat Equation as Repeated BVP}

\begin{thm}[Implicit Euler as BVP]\label{thm:heat-bvp}
The heat equation $u_t = u_{xx}$ with Dirichlet boundary conditions,
discretized via implicit Euler with time step $\Delta t$, is equivalent to
solving a BVP at each time step:
\[
(I - \Delta t \cdot d^2/dx^2) u^{n+1} = u^n.
\]
\end{thm}

\begin{proof}
The implicit Euler discretization of $u_t = u_{xx}$ is:
\[
\frac{u^{n+1} - u^n}{\Delta t} = u^{n+1}_{xx}
\quad\Longrightarrow\quad
u^{n+1} - \Delta t \cdot u^{n+1}_{xx} = u^n.
\]
Let $D = -d^2/dx^2$ with domain $\{u : u(0) = u(1) = 0\}$ and Green's
function $K = D^{-1}$ (the Fredholm kernel). Then:
\[
u^{n+1} + \frac{1}{\Delta t} \cdot K[u^{n+1}] = \frac{1}{\Delta t} \cdot
K[u^n].
\]
Setting $\alpha = 1/\Delta t$ and
$A[u] \equiv u + \alpha \cdot K[u] - \alpha \cdot K[u^n]$,
the equation $A[u^{n+1}] = 0$ is solved by
$\fix(T(A, \eta))$ with an appropriate $\eta$. Each time step is one
invocation of the BVP solver. The PDE adds nothing fundamentally new---it
applies the BVP operator $N_t$ times.
\end{proof}

% ═══════════════════════════════════════════════════════════════════════════
\section{Automatic Differentiation as a Term Operator}\label{sec:ad}
% ═══════════════════════════════════════════════════════════════════════════

\subsection{The Wengert Tape}

A \textbf{Wengert tape}~\cite{wengert1964} is a sequence of scalar
operations recorded during forward evaluation. The backward pass processes
the sequence in reverse, propagating partial derivatives via the chain rule.

Within the calculus, the tape is a mutable data structure attached to the
evaluation context. Each scalar operation that participates in the
computation records an entry:

\begin{defi}[Tape entry]
\begin{align*}
\mathsf{TapeEntry} ::= (&op, inputs, value, metadata) \\
\text{where } op \in \{&\mathsf{var}, \mathsf{add}, \mathsf{sub},
  \mathsf{mul}, \mathsf{div}, \mathsf{pow}, \mathsf{neg}, \\
  &\mathsf{sin}, \mathsf{cos}, \mathsf{exp}, \mathsf{log},
  \mathsf{sqrt}, \mathsf{tanh}, \mathsf{relu}\} \\
inputs &\subseteq \NN \quad\text{(tape indices of operands)} \\
value &\in \RR \quad\text{(cached forward value)} \\
metadata &\quad\text{(operand values for gradient rules)}
\end{align*}
\end{defi}

\begin{thm}[Reverse-mode correctness]\label{thm:tape-correct}
For any expression $e$ built from the tape-tracked operations, the reverse
pass computes the exact gradient $\nabla e$ with relative error bounded by
machine precision.
\end{thm}

\begin{proof}
Each operation $\mathit{op}$ has an associated backward rule that is the
exact partial derivative. For example:
\begin{align*}
\mathsf{add}(a, b):&\quad da \;+\!\!= d_{\mathit{out}},\;
  db \;+\!\!= d_{\mathit{out}} \\
\mathsf{mul}(a, b):&\quad da \;+\!\!= d_{\mathit{out}} \cdot b,\;
  db \;+\!\!= d_{\mathit{out}} \cdot a \\
\mathsf{tanh}(a):&\quad da \;+\!\!= d_{\mathit{out}} \cdot (1 -
  \mathsf{tanh}^2(a))
\end{align*}
By induction on the structure of the tape (which forms a DAG for
feedforward expressions), the reverse accumulation computes
$\partial e / \partial x_i$ for all inputs $x_i$ simultaneously. The error
is due solely to floating-point arithmetic. The full set of backward rules
is given in Appendix~\ref{sec:backward-rules}.
\end{proof}

\begin{cor}[Complexity]\label{cor:tape-complexity}
For an expression with $p$ scalar inputs and $n$ operations, the tape
computes all $p$ partial derivatives in $O(n)$ time---independent of $p$.
\end{cor}

\subsection{Grad as a Term}

\begin{defi}[Grad operator]
$\mathsf{Grad} : \Tcal \to \Tcal$ is the term constructor:
\[
\mathsf{Grad}(\mathsf{Lam}(\theta, body))
\]
which, when applied to a parameter vector $\theta_{\mathit{val}}$, returns
$\nabla_\theta\; body(\theta_{\mathit{val}})$.
\end{defi}

The semantics of $\mathsf{Grad}$ consists of three steps:
\begin{enumerate}
\item \textbf{Create} a fresh tape.
\item \textbf{Evaluate} $\mathit{body}$ with $\theta$ bound to tape-tracked
  variables, recording all scalar operations on the tape.
\item \textbf{Execute} the reverse pass, accumulating the gradient.
\end{enumerate}

\begin{thm}[Grad is a term operator]\label{thm:grad-term}
$\mathsf{Grad}$ is type-preserving: if $\mathit{body} : X \to \RR$, then
$\mathsf{Grad}(\mathsf{Lam}(\theta, \mathit{body})) : X \to X^*$. The
gradient is computed entirely within the term algebra without external
procedure calls.
\end{thm}

\begin{proof}
The forward pass records operations on the tape. The backward pass reads the
tape in reverse order. Both are implemented as reduction rules operating on
terms. No external procedure---no call to an AD library---is invoked. The
entire computation lives within the reduction system of
Section~\ref{sec:calculus}.
\end{proof}

The operational realization of $\mathsf{Grad}$ is a two-component system. A
$\mathsf{Tape}$ structure maintains a Wengert list of
$\mathsf{TapeEntry}$ objects. A $\mathsf{TapeVar}$ wraps a scalar value
with a tape reference and overloads the arithmetic operators ($+$, $-$, $*$,
$/$, exponentiation) and elementary functions ($\sin$, $\cos$, $\exp$,
$\tanh$, $\mathsf{relu}$). Each operation calls
\[
\mathsf{record}(op,\; result\_val,\; left\_id,\; right\_id,\; \dots)
\]
appending an entry to the tape. The forward pass builds the tape; the
backward pass replays it in reverse, seeding the last entry with gradient
$1.0$ and propagating via the chain rule (Appendix~\ref{sec:backward-rules}).

The key design choice is that $\mathsf{TapeVar}$ operations are syntactically
identical to ordinary arithmetic. An expression such as
$\theta_1 \cdot x + \theta_2$ evaluates to a scalar when $\mathsf{TapeVar}$
carries no tape reference (prediction mode), and builds the full computation
graph when it carries an active tape (training mode). This zero-overhead
switching enables a single forward function to serve both prediction and
gradient computation, without duplicating the model definition.

% ═══════════════════════════════════════════════════════════════════════════
\section{Deep Equilibrium Models as Nested Fixed Points}\label{sec:deq}
% ═══════════════════════════════════════════════════════════════════════════

\subsection{Forward Pass as Fix}

A Deep Equilibrium Model~\cite{bai2019} defines a layer implicitly:
\[
z^* = f(z^*, x, \theta).
\]
In the calculus, this is:
\[
z^* = \fix(\mathsf{Lam}(z,\; f(z, x, \theta))).
\]
The evaluation rule $\fix(f) \to \app{f}{\fix(f)}$ generates the sequence
$z_{k+1} = f(z_k, x, \theta)$. Convergence to a fixed point $z^*$ depends on
the contractivity of $f$ with respect to $z$.

\subsection{Gradient as Fix}

\begin{thm}[DEQ gradient via implicit function theorem]\label{thm:deq-grad}
The gradient of a loss $\mathcal{L}(z^*)$ with respect to parameters
$\theta$, where $z^* = \fix(\mathsf{Lam}(z,\; f(z, x, \theta)))$, is:
\[
\frac{\partial \mathcal{L}}{\partial \theta}
  = g^* \cdot \frac{\partial f}{\partial \theta}
\]
where the adjoint $g^*$ satisfies:
\[
g^* = \fix\!\left(\mathsf{Lam}\!\left(g,\;
  \frac{\partial \mathcal{L}}{\partial z^*} +
  g \cdot \frac{\partial f}{\partial z}\right)\right).
\]
\end{thm}

\begin{proof}
From $z^* = f(z^*, x, \theta)$, differentiate with respect to $\theta$:
\begin{align*}
\frac{\partial z^*}{\partial \theta}
  &= \frac{\partial f}{\partial z} \cdot \frac{\partial z^*}{\partial \theta}
   + \frac{\partial f}{\partial \theta} \\[4pt]
\Rightarrow\;
(I - \frac{\partial f}{\partial z}) \cdot \frac{\partial z^*}{\partial \theta}
  &= \frac{\partial f}{\partial \theta} \\[4pt]
\Rightarrow\;
\frac{\partial z^*}{\partial \theta}
  &= (I - \frac{\partial f}{\partial z})^{-1} \cdot
     \frac{\partial f}{\partial \theta}.
\end{align*}
Then:
\begin{align*}
\frac{\partial \mathcal{L}}{\partial \theta}
  &= \frac{\partial \mathcal{L}}{\partial z^*} \cdot
     \frac{\partial z^*}{\partial \theta} \\
  &= \frac{\partial \mathcal{L}}{\partial z^*} \cdot
     (I - \frac{\partial f}{\partial z})^{-1} \cdot
     \frac{\partial f}{\partial \theta} \\
  &= g^* \cdot \frac{\partial f}{\partial \theta},
\end{align*}
where $g^* = \frac{\partial \mathcal{L}}{\partial z^*} \cdot
(I - \frac{\partial f}{\partial z})^{-1}$.

Solving for $g^*$:
\begin{align*}
g^* \cdot (I - \frac{\partial f}{\partial z})
  &= \frac{\partial \mathcal{L}}{\partial z^*} \\
\Rightarrow\;
g^* &= \frac{\partial \mathcal{L}}{\partial z^*} +
      g^* \cdot \frac{\partial f}{\partial z}.
\end{align*}
This is a fixed-point equation in $g^*$. Hence
$g^* = \fix(\mathsf{Lam}(g,\;
\partial \mathcal{L}/\partial z^* + g \cdot \partial f/\partial z))$.
\end{proof}

\begin{cor}[Forward-backward unification]\label{cor:fb-unif}
For DEQ models, both the forward pass and the gradient computation are
instances of $\fix(T)$. The chain rule is replaced by another fixed point.
\end{cor}

\begin{thm}[Memory complexity]\label{thm:deq-memory}
The DEQ gradient computation stores only the final fixed point $z^*$, not
the intermediate iterates $z_1, \dots, z_{K-1}$. The memory complexity is
$O(1)$ in the number of forward iterations.
\end{thm}

\begin{proof}
The adjoint equation depends on $f$ only through $\partial f/\partial z$
evaluated at $z^*$. The intermediate values $z_k$ are not needed. The
$\fix$ expansion for the gradient produces its own iterates $g_k$, but
these do not require storing $z_1, \dots, z_{K-1}$.
\end{proof}

The nested $\fix$ structure reveals a fundamental equivalence: training a
DEQ model consists of two nested invocations of the T-transform iteration.
The outer loop trains parameters ($\theta \leftarrow \theta - \eta \cdot
\nabla_\theta \mathcal{L}$), and the inner loop computes the adjoint
($g_{k+1} = \partial \mathcal{L}/\partial z^* + g_k \cdot \partial
f/\partial z$). Both are $\fix(T)$ for different choices of $A$. The
original DEQ formulation~\cite{bai2019} requires solving a linear system
$(I - \partial f/\partial z)^{-1}$ or using an iterative solver; in our
calculus, the linear solve \emph{is} the $\fix$ expansion---no distinction
between ``direct'' and ``iterative'' solvers exists. The solver is the
primitive.

% ═══════════════════════════════════════════════════════════════════════════
\section{Implementation}\label{sec:implementation}
% ═══════════════════════════════════════════════════════════════════════════

The entire formal system is implemented in approximately 1,500 lines of
code with a single numerical dependency. The T-transform and its fixed-point
solver together require 29 lines.

\subsection{From Theory to Computation}

The computational architecture mirrors the theoretical structure: term
constructors (Section~2.1) become datatype definitions, reduction rules
(Section~2.2) become an evaluator, and the T-transform
(Section~\ref{sec:ttransform}) compiles to a kernel of under 30 lines.

\subsection{Differentiation Strategies}

The system provides two gradient-computation strategies, both conforming to
the same abstract interface:

\begin{defi}[Model interface]
A \textbf{model} is a triple $(\theta_0, f, \nabla f)$ where
$\theta_0 \in \RR^p$ is the initial parameter vector,
$f : \RR^p \times \RR \to \RR$ is the forward function, and
$\nabla f : \RR^p \to \RR^p$ computes the gradient of a loss functional
with respect to $\theta$.
\end{defi}

\textbf{Finite-difference gradient.} For any forward function $f(\theta, x)$,
the gradient is approximated by central differences:
\[
(\nabla_\theta \mathcal{L})_i \approx
\frac{\mathcal{L}(\theta + h e_i) - \mathcal{L}(\theta - h e_i)}{2h},
\]
where $e_i$ is the $i$-th standard basis vector and $h \approx 10^{-5}$.
Complexity is $O(p)$ per gradient evaluation with $O(h^2)$ accuracy. No
tracing, no tape, no special syntax is required---any computable function
$(\theta, x) \mapsto y$ can be differentiated.

\textbf{Tape-based gradient.} For a forward function expressed in
$\mathsf{TapeVar}$ arithmetic (Section~\ref{sec:ad}), all scalar operations
are recorded on a Wengert tape during one forward pass. A single reverse
sweep---$\mathsf{backward}()$---computes all $p$ partial derivatives
simultaneously via the chain rule. Complexity is $O(n)$ where $n$ is the
number of scalar operations in the forward pass, independent of $p$. For a
typical neural network layer with weight matrix $W \in \RR^{m \times n}$ and
bias $b \in \RR^n$, the forward pass performs $\approx 2mn + n$ scalar
operations, and the reverse pass costs roughly the same---the ``cheap
gradient principle''~\cite{griewank2008}.

Both strategies are instances of the same $\mathsf{Grad}$ operator
(Definition~5.2), differing only in how they evaluate the derivative.
The training algorithm is agnostic to this choice.

\subsection{Training as T-Transform Iteration}

Training a model reduces to the T-transform iteration with
$A = \nabla_\theta \mathcal{L}$:
\begin{tabbing}
\quad \= \textbf{for} \= $t = 0, 1, \dots, N-1$: \\
\> \> $g_t \leftarrow \nabla_\theta \mathcal{L}(\theta_t)$ \\
\> \> $\theta_{t+1} \leftarrow \theta_t - \eta \cdot g_t$
\end{tabbing}
The parameter $\eta$ is the learning rate, time step, and damping factor
simultaneously (Theorem~\ref{thm:eta}). Each iteration is one application of
$T_{\nabla \mathcal{L}, \eta}$. The algorithm does not inspect whether $g_t$
came from finite differences or a reverse-mode tape---the gradient is a
vector in $\RR^p$, and the update is vector subtraction. The model provides
$g_t$; the trainer provides $\eta$ and the loop; the T-transform provides
the guarantee that $\fix(T_{\nabla \mathcal{L}, \eta})$ is a stationary
point of $\mathcal{L}$.

\subsection{PDE Solving}

\begin{center}
\begin{tabular}{l c c c}
\toprule
Equation & $A$ operator & Iterations & Error \\
\midrule
$x^3 - 2x - 5 = 0$ &
  $\mathsf{Lam}(x,\; x^3 - 2x - 5)$ & 9 & $f(x^*) = 2.1\times 10^{-9}$ \\
$y' = y,\; y(0) = 1$ &
  Picard residual & 11 & $L^2$ error $2.0\times 10^{-5}$ \\
$-u'' = 1,\; u(0) = u(1) = 0$ &
  BVP residual & 24 & $L^2$ error $5.4\times 10^{-9}$ \\
Heat eq.\ (3 steps) &
  Implicit Euler residual & 5/step & $O(\Delta t)$ error \\
\bottomrule
\end{tabular}
\end{center}

\subsection{Neural Network Training}

\begin{center}
\begin{tabular}{l c c c}
\toprule
Task & Model & Parameters & Gradient method \\
\midrule
Linear regression $(y = 2x + 1)$ &
  FunctionModel & 2 & Finite diff.\ \\
$\sin(x)$ regression &
  MLP[1,16,1] & 49 & Tape ($O(1)$) \\
Binary classification &
  Linear + BCE & 2 & Finite diff.\ \\
DEQ $\tanh$ fit &
  DEQ function & 3 & Finite diff.\ \\
Gradient accuracy check &
  TapeModel & 49 & Tape vs.\ FD \\
CIFAR-10 &
  MLP[256,64,10] & 17,162 & Tape ($O(1)$) \\
\bottomrule
\end{tabular}
\end{center}

\begin{center}
\begin{tabular}{l c}
\toprule
Result & Value \\
\midrule
Linear regression: $\hat{a}, \hat{b}$ & 1.94, 1.08 \\
$\sin(x)$ regression: final MSE & 0.061 \\
Binary classification: accuracy & 93.3\% \\
DEQ $\tanh$ fit: final MSE & 0.009 \\
Tape vs.\ finite difference: relative error & $\sim 10^{-11}$ \\
\bottomrule
\end{tabular}
\end{center}

The tape-based reverse-mode AD computes exact gradients (relative error
$\sim 10^{-11}$ versus finite differences) in $O(1)$ time in the number of
parameters, independent of the forward expression size.

\subsection{Reproducibility}

The complete implementation, including all experiments, is available as
ancillary material accompanying this submission. All numerical results
reported in this section are reproducible from the provided code.

% ═══════════════════════════════════════════════════════════════════════════
\section{Discussion}\label{sec:discussion}
% ═══════════════════════════════════════════════════════════════════════════

\subsection{Related Formal Systems}

The presentation of $\fix$ as a primitive constructor differs from its
standard treatment in domain theory~\cite{scott1976,plotkin1977}, where the
least fixed-point operator $\mathbf{Y}$ is definable or axiomatized as
$\mathbf{Y}(f) = f(\mathbf{Y}(f))$ in a typed setting. Our approach is
closer to the spirit of the untyped $\lambda$-calculus, where recursion is
freely available, but we retain the explicit $\fix$ constructor to make the
fixed-point structure visible to the syntax.

The $(D, C)$ decomposition echoes the formal theory of boundary value
problems in functional analysis~\cite{stakgold1998}, but we internalize the
decomposition within the term calculus as an operator factorization
(Lemma~\ref{lem:op-factor}): $D$ and $C$ are not merely a pair of operators
but factors $D_0 + BC$ of a single bounded linear operator on a normed space.

The present paper provides an operational semantics (reduction rules) and a
denotational semantics restricted to complete metric spaces (Banach fixed-point
interpretation, Remark~\ref{rem:fix-semantics}). A full denotational semantics
in the sense of domain theory or traced monoidal categories---with soundness and
completeness relative to the operational semantics---remains for future work
(Section~\ref{sec:future}).

\subsection{Future Work}\label{sec:future}

\begin{enumerate}
\item \textbf{Denotational semantics.} The present paper provides an
  operational semantics (reduction rules) and a Banach fixed-point
  interpretation of $\fix$ (Remark~\ref{rem:fix-semantics}). A full
  denotational semantics---with soundness and completeness theorems relating
  the reduction system to a domain-theoretic or traced monoidal category
  model---is the principal next step. The T-transform $T[u] = u - \eta \cdot
  A[u]$ is a natural transformation in the category of normed vector spaces,
  and $\fix$ is a trace in the sense of Joyal, Street, and
  Verity~\cite{joyal1996}, suggesting that traced monoidal categories and the
  geometry of interaction provide the right semantic foundation.

\item \textbf{Type theory.} The current system is unityped. Adding a type
  system with $\fix$ as a type constructor ($\mu$-types) would connect to
  domain theory and inductive/coinductive types~\cite{pierce2002,harper2016}.

\item \textbf{Convergence theory.} The formal correctness of $\fix(T)$
  (Theorem~\ref{thm:fix-solves}) is independent of convergence. A complete
  convergence theory would characterize the class of operators $A$ for which
  $T_{A,\eta}$ is a contraction, and the optimal $\eta$ schedules for
  non-convex operators.

\item \textbf{Compiler.} The term evaluator is an interpreter. A compiler
  from Term expressions to optimized numerical code (JIT or ahead-of-time)
  would eliminate the interpreter overhead for deployed models.

\item \textbf{Probabilistic extension.} Adding a noise term
  $T[u] = u - \eta \cdot A[u] + \sqrt{2\eta} \cdot \xi$ yields Langevin
  dynamics. $\fix(T)$ becomes the stationary distribution---Bayesian
  inference as fixed point.

\item \textbf{Higher-order operators.} The current calculus applies
  $\mathsf{Integ}$ and $\mathsf{Diff}$ to first-order terms. Extending to
  higher-order operators (biharmonic, polyharmonic) follows the same
  template: each operator is a Term whose application dispatches to the
  corresponding Green's kernel.
\end{enumerate}

\subsection{Comparison with Machine Learning}

While the contributions of this paper are formal, they bear directly on
computational practice. The standard neural network training loop---forward
pass, loss, backward pass, parameter update---is exactly the T-transform
iteration $u_{t+1} = u_t - \eta \cdot A[u_t]$ with $A = \nabla_\theta
\mathcal{L}$. That this is not a coincidence but a formal identity means that
any improvement to fixed-point convergence (better $\eta$ schedules,
preconditioners, acceleration methods like Anderson mixing~\cite{anderson1965})
immediately improves training across all model architectures, without
modifying the model.

% ═══════════════════════════════════════════════════════════════════════════
\section{Conclusion}
% ═══════════════════════════════════════════════════════════════════════════

We have defined a term calculus in which the fixed-point combinator $\fix$
is a primitive syntactic constructor, the T-transform
$T[u] = u - \eta \cdot A[u]$ is the canonical affine residual operator
embedding operator equations into discrete dynamical systems, and the single
construct $\fix(T)$ solves scalar equations, initial value problems, boundary
value problems, PDEs, and the gradient computations of deep equilibrium
models---all as instances of the same operator flow.

The central insight is that \textbf{the boundary between ``solving an
equation'' and ``training a model'' is not a mathematical distinction but an
artifact of our computational tools.} What is unified is not the structure of
the equations but the \emph{evolution in function space} that solves them. When
$\fix$ is a primitive of the language, solving \emph{is} computing.

% ── Acknowledgments ──
\subsection*{Acknowledgments}
The author thanks the LMCS editors and anonymous reviewers for their
constructive feedback.

% ── Appendix ──
\appendix
\section{Backward Rules for the Wengert Tape}\label{sec:backward-rules}

The following rules are applied during the reverse pass of the tape. Each
rule receives the output gradient $d_{\mathit{out}}$ and accumulates into
the gradients of the inputs.

\begin{center}
\begin{tabular}{l l}
\toprule
Operation & Backward rule \\
\midrule
$\mathsf{add}(a, b)$ &
  $\mathit{da} \;+\!\!= d_{\mathit{out}};\;
   \mathit{db} \;+\!\!= d_{\mathit{out}}$ \\[3pt]
$\mathsf{sub}(a, b)$ &
  $\mathit{da} \;+\!\!= d_{\mathit{out}};\;
   \mathit{db} \;-\!\!= d_{\mathit{out}}$ \\[3pt]
$\mathsf{mul}(a, b)$ &
  $\mathit{da} \;+\!\!= d_{\mathit{out}} \cdot b;\;
   \mathit{db} \;+\!\!= d_{\mathit{out}} \cdot a$ \\[3pt]
$\mathsf{div}(a, b)$ &
  $\mathit{da} \;+\!\!= d_{\mathit{out}} / b;\;
   \mathit{db} \;-\!\!= d_{\mathit{out}} \cdot a / b^2$ \\[3pt]
$\mathsf{pow}(a, b)$ &
  $\mathit{da} \;+\!\!= d_{\mathit{out}} \cdot b \cdot a^{b-1}$ ($a \neq 0$);
  \\
& $\mathit{db} \;+\!\!= d_{\mathit{out}} \cdot a^b \cdot \ln a$ ($a > 0$) \\[3pt]
$\mathsf{neg}(a)$ &
  $\mathit{da} \;-\!\!= d_{\mathit{out}}$ \\[3pt]
$\mathsf{sin}(a)$ &
  $\mathit{da} \;+\!\!= d_{\mathit{out}} \cdot \cos a$ \\[3pt]
$\mathsf{cos}(a)$ &
  $\mathit{da} \;-\!\!= d_{\mathit{out}} \cdot \sin a$ \\[3pt]
$\mathsf{exp}(a)$ &
  $\mathit{da} \;+\!\!= d_{\mathit{out}} \cdot e^a$ \\[3pt]
$\mathsf{log}(a)$ &
  $\mathit{da} \;+\!\!= d_{\mathit{out}} / a$ \\[3pt]
$\mathsf{sqrt}(a)$ &
  $\mathit{da} \;+\!\!= d_{\mathit{out}} / (2\sqrt{a})$ \\[3pt]
$\mathsf{tanh}(a)$ &
  $\mathit{da} \;+\!\!= d_{\mathit{out}} \cdot (1 - \tanh^2 a)$ \\[3pt]
$\mathsf{relu}(a)$ &
  $\mathit{da} \;+\!\!= d_{\mathit{out}}$ if $a > 0$, else $0$ \\
\bottomrule
\end{tabular}
\end{center}

The backward rules are the exact partial derivatives. The only source of
error in the computed gradient is floating-point arithmetic.

\bibliographystyle{alphaurl}
\bibliography{paper}

\end{document}
